\section{Historia y cambios}\label{sec:cambios}

\subsection{De lenguaje formal a escritura para apuntes}
La primera versión de KalaScript (etiqueta \texttt{v1-formal} en git) era un
fragmento de lógica formal: cuantificadores, lógica modal con modelos de
Kripke, programas con \emph{mientras} y un catálogo de relaciones numeradas
R1--R9. Estaba bien construida, pero no servía para tomar apuntes: redibujaba
\(\forall\), \(\exists\), \(\neg\) y \(\to\), y cada relación nueva exigía un
número y una definición formal.

En paralelo existió otra versión (en la memoria USB), con una sintaxis textual
tipo grafo, respaldo fonético y un proyecto de corpus. De ella tomé tres ideas:
distinguir de dónde viene una afirmación y qué tan segura es, pensar en
sucesos con participantes (la plantilla de suceso) y que el sistema crezca a
partir de apuntes reales.

\subsection{Qué se conservó}
\begin{itemize}
\item Trece signos con su dibujo intacto: K01--K10, K16, K19 y K21. Siguen
 siendo las formas B de Colección (K06) y Fuente (K16) que elegí al practicar.
\item Las distinciones de fondo: negar no es oponer, sucesión no es causa, parte
 no es miembro, fuente no es certeza.
\item La cuadrícula de \(64\times64\), el registro único de signos y los
 dibujos en TikZ.
\end{itemize}

\subsection{Qué cambió}
\begin{itemize}
\item \textbf{Nombres más cortos} para cinco signos, sin cambiar su idea:
 Entidad pasa a llamarse Cosa (K01), Colección pasa a Lista (K06), Espacio
 pasa a Lugar (K08), Referencia pasa a Ver (K10) y Fuente pasa a Según (K16).
\item \textbf{Usos ampliados:} Relación (K03) ya no exige un número de
 catálogo; Secuencia (K21), que ordenaba programas, pasa a ser Luego y ordena
 cualquier sucesión; Pregunta (K19) cubre también dudas y cosas no entendidas.
\item \textbf{Retirados:} K11--K15, K17, K18, K20 y K22 (anexo~\ref{sec:retirados}).
\item \textbf{Nuevos:} trece signos, K23--K35: siete conectores, una categoría
 y cinco marcas de margen.
\item \textbf{Nuevas reglas:} hoja con margen, símbolos matemáticos con
 significado fijo, abreviaturas del español, plantillas y método para añadir
 signos.
\end{itemize}

\subsection{Inspiraciones}
La taquigrafía Gregg y Teeline (escribir rápido), los Blissymbolics (signos
por significado), los grafos existenciales de Peirce (el alcance como algo
visible), el método Cornell (margen para repasar), el \emph{bullet journal}
(marcas de margen para tareas) y la notación de Iverson (pocos signos bien
elegidos).
